% ============================================================
% DOCUMENTCLASS
% ============================================================
\documentclass[12pt,oneside]{book}

% ============================================================
% METADATOS
% ============================================================
\newcommand{\universidad}{Universidad de Buenos Aires (UBA)}
\newcommand{\facultad}{Facultad de Derecho}
\newcommand{\carrera}{Carrera de Abogacía --- Módulo de Derecho Público}
\newcommand{\materia}{Derecho Constitucional Profundizado}
\newcommand{\titulo}{El Recurso Extraordinario Federal y la Cuestión Federal}
\newcommand{\autor}{Isaac Edgar Camacho Ocampo}
\newcommand{\anio}{2026}

% ============================================================
% PAQUETES
% ============================================================
\usepackage[utf8]{inputenc}
\usepackage[T1]{fontenc}
\usepackage[spanish,es-nodecimaldot]{babel}

\usepackage[
    top=2.5cm,
    bottom=2.5cm,
    left=3cm,
    right=2.5cm,
    headheight=15pt
]{geometry}

\usepackage{amsmath}
\usepackage{amssymb}
\usepackage{mathtools}

\usepackage{graphicx}
\usepackage{float}
\usepackage{caption}
\usepackage{subcaption}

\usepackage{booktabs}
\usepackage{array}
\usepackage{longtable}
\usepackage{tabularx}

\usepackage{enumitem}

\usepackage{xcolor}
\usepackage[most]{tcolorbox}

\usepackage{fancyhdr}

\usepackage{ragged2e}

\usepackage[
    colorlinks=true,
    linkcolor=blue!50!black,
    citecolor=green!40!black,
    urlcolor=blue!60!black,
    pdftitle={\titulo},
    pdfauthor={\autor},
    pdfsubject={Apuntes universitarios},
    bookmarks=true
]{hyperref}

% ============================================================
% CONFIGURACIÓN
% ============================================================
\graphicspath{
    {./img/}
    {./graficos/}
    {./figuras/}
}

\setcounter{tocdepth}{2}

\setlength{\parindent}{0pt}
\setlength{\parskip}{0.5em}

\setlist[itemize]{
    topsep=0.3em,
    itemsep=0.2em
}

\setlist[enumerate]{
    topsep=0.3em,
    itemsep=0.2em
}

\clubpenalty=10000
\widowpenalty=10000

% ============================================================
% COMANDOS
% ============================================================
\newcommand{\abs}[1]{\left|#1\right|}
\newcommand{\norm}[1]{\left\lVert#1\right\rVert}
\newcommand{\vect}[1]{\mathbf{#1}}
\newcommand{\deriv}[2]{\frac{d#1}{d#2}}
\newcommand{\pderiv}[2]{\frac{\partial#1}{\partial#2}}

% ============================================================
% ENTORNOS / CAJAS
% ============================================================
\newtcolorbox{definition}[1]{
    enhanced,
    breakable,
    title={Definición: #1},
    fonttitle=\bfseries,
    colback=blue!3,
    colframe=blue!50!black,
    boxrule=0.8pt,
    arc=2mm
}

\newtcolorbox{important}{
    enhanced,
    breakable,
    title={Importante},
    fonttitle=\bfseries,
    colback=yellow!8,
    colframe=orange!70!black,
    boxrule=0.8pt,
    arc=2mm
}

\newtcolorbox{example}[1]{
    enhanced,
    breakable,
    title={Ejemplo: #1},
    fonttitle=\bfseries,
    colback=green!4,
    colframe=green!40!black,
    boxrule=0.8pt,
    arc=2mm
}

\newtcolorbox{formula}[1]{
    enhanced,
    breakable,
    title={Fórmula / Regla: #1},
    fonttitle=\bfseries,
    colback=purple!4,
    colframe=purple!50!black,
    boxrule=0.8pt,
    arc=2mm
}

\newtcolorbox{exercise}[1]{
    enhanced,
    breakable,
    title={Ejercicio: #1},
    fonttitle=\bfseries,
    colback=gray!5,
    colframe=gray!60!black,
    boxrule=0.8pt,
    arc=2mm
}

\newtcolorbox{note}{
    enhanced,
    breakable,
    title={Nota},
    fonttitle=\bfseries,
    colback=white,
    colframe=gray!50,
    boxrule=0.6pt,
    arc=2mm
}

% ============================================================
% CONFIGURACIÓN FINAL (encabezados y pies)
% ============================================================
\pagestyle{fancy}
\fancyhf{}

\fancyhead[L]{\small\nouppercase{\leftmark}}
\fancyhead[R]{\small\materia}

\fancyfoot[C]{\thepage}

\renewcommand{\headrulewidth}{0.4pt}
\renewcommand{\footrulewidth}{0pt}

\fancypagestyle{plain}{
    \fancyhf{}
    \fancyfoot[C]{\thepage}
    \renewcommand{\headrulewidth}{0pt}
}

% ============================================================
% DOCUMENT
% ============================================================
\begin{document}

% ---------------- PORTADA ----------------
\begin{titlepage}

\thispagestyle{empty}

\begin{center}

\vspace*{1cm}

{\large \textsc{\universidad}}

\vspace{0.2cm}

{\large \textsc{\facultad}}

\vspace{0.2cm}

{\large \carrera}

\vspace{3cm}

{\Huge \textbf{\materia}}

\vspace{1cm}

{\LARGE \titulo}

\vspace{0.8cm}

\textit{Comprender los conceptos es el primer paso para convertir el estudio en conocimiento.}

\vspace{1.2cm}

\rule{0.70\textwidth}{0.5pt}

\vspace{0.5cm}

{\large Apuntes de estudio}

\vspace{0.5cm}

\rule{0.70\textwidth}{0.5pt}

\vfill

{\large \autor}

\vspace{0.3cm}

{\large \anio}

\end{center}

\vfill

\begin{flushleft}
\textit{Este es un modesto aporte a los estudiantes de ``Derecho Constitucional Profundizado'' no pretende sustituir el material oficial de la materia.}
\end{flushleft}

\end{titlepage}

% ---------------- ÍNDICE ----------------
\tableofcontents

% ============================================================
% CONTENIDO
% ============================================================

\chapter{Marco teórico y metodológico}

\section{Presentación de la unidad}

La presente unidad reúne el análisis de cinco casos prácticos expuestos en una clase de exposición grupal, todos ellos vinculados a un mismo eje conceptual: la noción de \textbf{cuestión federal} y su canalización procesal a través del \textbf{recurso extraordinario federal}. Los temas específicos abordados son los siguientes:

\begin{itemize}
    \item El sistema de control de constitucionalidad argentino y los requisitos del recurso extraordinario federal (art. 14, Ley 48).
    \item La noción de cuestión federal simple, compleja directa y compleja indirecta, a partir de cinco casos reales analizados en clase.
    \item El derecho a la vivienda digna y los grupos prioritarios en la Ciudad Autónoma de Buenos Aires.
    \item Los decretos de necesidad y urgencia y sus límites materiales (art. 99 inc. 3 CN), aplicados a la ciudadanía y las migraciones.
    \item La movilidad jubilatoria y el control de constitucionalidad de normas previsionales (caso ``Badaro'').
    \item El poder tributario municipal, la cláusula comercial y la prohibición de aduanas interiores.
    \item La autonomía universitaria, la insistencia legislativa (art. 83 CN) y las medidas cautelares frente al Poder Ejecutivo.
\end{itemize}

\begin{note}
\textbf{Relación entre los temas de la clase.} Los cinco casos analizados pueden entenderse como manifestaciones distintas de un mismo conflicto estructural: una norma de rango inferior --- una ordenanza, un decreto o una ley reglamentaria --- que entra en tensión con una norma de rango superior --- la Constitución Nacional, un tratado con jerarquía constitucional o una ley federal. En todos los casos aparece la misma pregunta (¿esta norma inferior contradice la norma superior?) y la misma vía de solución: la cuestión federal y el recurso extraordinario federal ante la Corte Suprema de Justicia de la Nación.
\end{note}

\begin{important}
\textbf{Relevancia profesional de la unidad.} El valor de esta unidad no se limita a la memorización de artículos: cada uno de los cinco casos analizados --- vivienda, migraciones, jubilaciones, tributos municipales y financiamiento universitario --- representa un conflicto real entre una norma inferior y la Constitución. Reconocer esa tensión, y saber canalizarla mediante el recurso extraordinario federal, constituye una destreza exigida en la práctica profesional de la abogacía, tanto en el ejercicio privado como en el sector público.
\end{important}

\section{Encuadre metodológico de la clase}

La clase se desarrolla bajo la modalidad de exposición grupal de casos reales. El docente indica que cada grupo debe identificar, en el caso elegido, si existe una cuestión federal y de qué tipo --- simple, compleja directa o compleja indirecta ---, describir el recorrido procesal que llevó (o que podría llevar) el caso hasta la Corte Suprema de Justicia de la Nación, y fundamentar la eventual procedencia del recurso extraordinario federal. A lo largo del encuentro se suceden cinco exposiciones grupales, cada una referida a un caso distinto, con intervenciones y correcciones del docente luego de cada presentación.

\section{La cuestión federal y el recurso extraordinario federal: nociones generales}

\begin{definition}{Cuestión federal}
Es el conflicto entre una norma de rango inferior (ordenanza, decreto, ley local o federal) y la Constitución Nacional, un tratado o una ley federal, cuya resolución habilita, eventualmente, el recurso extraordinario federal ante la Corte Suprema de Justicia de la Nación.
\end{definition}

\begin{definition}{Tipos de cuestión federal}
Los apuntes distinguen dos tipos de cuestión federal que se desarrollan efectivamente a lo largo de los casos analizados:
\begin{itemize}
    \item \textbf{Compleja directa}: existe una colisión inmediata entre una norma y la Constitución Nacional, sin normas intermedias. Se verifica, según lo expuesto en clase, en los casos de vivienda, ciudadanía, jubilaciones y universidades.
    \item \textbf{Compleja indirecta}: existe una colisión entre dos normas infraconstitucionales de distinta jerarquía, como ocurre entre la ordenanza municipal y la ley de coparticipación en el caso de la tasa vial.
\end{itemize}
% TODO: contenido incompleto o ambiguo en las notas originales. El apunte menciona también la categoría de "cuestión federal simple" entre los temas tratados, pero no desarrolla su definición ni la aplica a ninguno de los cinco casos expuestos.
\end{definition}

\begin{formula}{Requisitos del recurso extraordinario federal (art. 14, Ley 48)}
Conforme a lo señalado en los apuntes, la procedencia del recurso extraordinario federal exige:
\begin{itemize}
    \item la existencia de un caso judicial;
    \item una sentencia definitiva dictada por el superior tribunal de la causa;
    \item el agotamiento de las instancias ordinarias previas; y
    \item una relación directa entre la cuestión federal planteada y la resolución del caso.
\end{itemize}
\end{formula}

% ============================================================
\chapter[Caso 1 --- Vivienda digna en la CABA]{Caso 1 --- El derecho a la vivienda digna y los grupos prioritarios en la Ciudad de Buenos Aires}

El primer grupo expositor presenta un caso vinculado a su experiencia laboral en el Poder Judicial de la Ciudad Autónoma de Buenos Aires, en el ámbito del Ministerio Público de la Defensa. Por razones de discrecionalidad profesional, no se identifican los nombres reales de las partes involucradas.

\section{Los hechos}

La parte actora es una mujer de origen ruso, nacionalizada argentina, que en octubre de 2017 ingresó al país como peticionante de refugio, a raíz de la persecución social y el abuso policial que sufría en su país de origen por su orientación sexual y su participación en organizaciones LGBT. En el año 2020 inició un amparo habitacional, momento en el que interviene el Ministerio Público de la Defensa. Recientemente le fue expedido un certificado de discapacidad, por padecer osteonecrosis y coxartrosis --- artrosis de cadera ---, lo que le genera anormalidades en la marcha y la movilidad y le impide trabajar.

\section{El recorrido procesal}

En primera instancia del fuero contencioso administrativo y tributario de la Ciudad no se le reconoce el beneficio. Apelada esa decisión por el Ministerio Público de la Defensa, la Sala I de la Cámara hace lugar al reclamo: condena al Gobierno de la Ciudad de Buenos Aires a presentar ante el juzgado de origen una propuesta para hacer frente a la obligación de brindar alojamiento a la amparista, equiparando su situación a la de los grupos prioritarios reconocidos por la Ley 4036 --- mayores de 60 años, personas con discapacidad y, por vía jurisprudencial, mujeres en situación de violencia de género ---.

El Gobierno de la Ciudad interpone un recurso de inconstitucionalidad contra esa resolución, que es denegado; ante esa denegatoria, presenta un recurso de queja ante el Tribunal Superior de Justicia. El Tribunal Superior admite la queja, revoca la sentencia de Cámara y devuelve las actuaciones para que, por intermedio de otros jueces, se dicte un nuevo fallo. En consecuencia, el Ministerio Público de la Defensa presenta el recurso extraordinario federal, ya en agosto de 2026.

\section{Las normas puestas en juego}

\begin{formula}{Artículo 14 bis, Constitución Nacional (fragmento invocado)}
``El Estado otorgará los beneficios de la seguridad social, que tendrá carácter de integral e irrenunciable [\ldots]; la protección integral de la familia [\ldots] y el acceso a una vivienda digna.''

Citado junto con el art.\ 16 CN (igualdad ante la ley) como fundamento de la posición de la actora.
\end{formula}

\begin{formula}{Ley 4036 --- Protección Integral de los Derechos Sociales de la Ciudad de Buenos Aires}
Según lo referido en la exposición, reconoce como grupos prioritarios de atención a las personas mayores de 60 años y a las personas con discapacidad; por vía de un fallo posterior se incorporó también a las mujeres en situación de violencia de género. Es una norma de derecho público local que estructura el sistema de asistencia social porteño.
\end{formula}

\begin{formula}{Arts. 17 y 18, Constitución de la Ciudad Autónoma de Buenos Aires}
Según lo señalado en el recurso extraordinario federal, estos artículos establecen que el Gobierno de la Ciudad debe atender a las situaciones de diferencia social y garantizar el acceso prioritario a las políticas sociales a las personas en situación de vulnerabilidad, con obligaciones de cuidado y asistencia a su cargo. Se trata de la constitución local invocada como fuente directa de la obligación estatal en materia habitacional.
\end{formula}

\begin{formula}{Pacto Internacional de Derechos Económicos, Sociales y Culturales (PIDESC) y CEDAW}
El PIDESC consagra la obligación de los Estados parte de garantizar el goce de los derechos económicos, sociales y culturales --- entre ellos, el acceso a una vivienda digna --- ``hasta el más alto grado de satisfacción posible'', incluso frente a limitaciones presupuestarias. La CEDAW (Convención para la Eliminación de Todas las Formas de Discriminación contra la Mujer) fue invocada por la condición de la actora como mujer perteneciente al colectivo LGBT. Son tratados con jerarquía constitucional en los términos del art.\ 75 inciso 22 de la Constitución Nacional.
% TODO: contenido incompleto o ambiguo en las notas originales. El apunte original consigna, en un pasaje, la sigla "CEDAU" para el mismo tratado citado como "CEDAW" en el título de esta sección.
\end{formula}

\section{El debate en clase}

El argumento central del recurso sostiene que la Constitución no establece ningún ranking ni orden de prelación para acceder al beneficio: ni el art.\ 14 bis, ni el art.\ 16, ni los tratados internacionales invocados establecen un orden de prioridad entre las personas vulnerables.

\begin{important}
\textbf{Criterio hermenéutico aplicado.} Se explicó que existe un criterio de interpretación de la norma según el cual, allí donde la ley no establece distinciones, no corresponde introducirlas por vía interpretativa. En otras palabras, si la ley no prevé sujetos que accedan prioritariamente a un derecho, no puede invocarse esa ausencia de previsión como fundamento para denegarlo.
\end{important}

\begin{note}
\textbf{El PIDESC aplicado al caso.} Tratándose de derechos de amplio espectro íntimamente ligados al gasto estatal, existe una obligación internacional de los Estados de orientar sus políticas hacia el ``más alto grado de satisfacción'' de esos derechos. Siendo la Argentina parte del PIDESC, sus estados locales --- como la Ciudad de Buenos Aires --- también deben sostener ese principio, aun cuando existan limitaciones financieras: lo que no puede faltar es una política orientada en ese sentido.
\end{note}

Respecto de la condición migratoria de la actora, se aclaró un matiz relevante del caso: la actora ya se encuentra nacionalizada en Argentina. Presentó el amparo habitacional en 2020, pero recién en 2026 le fue otorgado el certificado de discapacidad; por ese motivo no se la admitía dentro del grupo prioritario. La denegación inicial no obedeció a su condición de nacida en el extranjero, sino a que en su momento no contaba con dicho certificado.

\section{Cierre del caso}

\begin{important}
\textbf{Síntesis.} Se trata de derecho público local: la interpretación de la Constitución de la Ciudad y de sus deberes en materia de vivienda digna. La causa tuvo una primera instancia denegatoria, fue admitida en Cámara y luego revocada por el Tribunal Superior de Justicia --- máximo tribunal local en materia de derecho administrativo ---, con lo cual se encuentra reunido el requisito de sentencia definitiva de última instancia para habilitar el recurso extraordinario federal, en trámite ante ese mismo Tribunal Superior.
\end{important}

% ============================================================
\chapter[Caso 2 --- Ciudadanía y migraciones]{Caso 2 --- Ciudadanía, migraciones y decretos de necesidad y urgencia}

\section{Los hechos}

El segundo grupo expone el caso de un hombre de nacionalidad china que ingresó al país en 2015 y se radicó en la provincia de Entre Ríos, donde abrió un supermercado, ejerció el comercio, tributó regularmente y formó una familia, desarrollando una vida normal durante varios años. Tras cumplir el plazo mínimo legal de residencia continua, solicitó formalmente la carta de ciudadanía ante la Justicia Federal, que históricamente era la competente para otorgarla.

El conflicto se origina cuando, en 2025, el Poder Ejecutivo dicta el Decreto de Necesidad y Urgencia (DNU) N.\textsuperscript{o} 366/2025, que transfiere la competencia para tramitar y otorgar la ciudadanía y la naturalización --- históricamente a cargo del Poder Judicial --- a la órbita del Poder Ejecutivo, concretamente a la Dirección Nacional de Migraciones, un órgano de naturaleza administrativa.

\section{El recorrido procesal}

Tramitada la solicitud por vía administrativa y agotada esa instancia, se le deniega la carta de ciudadanía por haber ingresado al país por un ``paso de los libres'', sin registro migratorio de ingreso.

\begin{formula}{Art. 29, Ley 25.871 --- Ley General de Migraciones}
Según lo referido en la exposición grupal, la norma dispone que no puede permanecer en el territorio nacional quien haya ingresado eludiendo el control migratorio, tenga una prohibición de ingreso vigente o presente documentación falsa o adulterada. Fue invocada por la autoridad migratoria como fundamento de la denegatoria de la ciudadanía.
\end{formula}

Al no ejecutarse la orden de expulsión, el interesado habilita la vía judicial federal y se presenta ante el Tribunal Federal N.\textsuperscript{o} 1 de Entre Ríos, aclarando que nunca fue efectivamente expulsado del país. Aun así, la primera instancia vuelve a rechazar el pedido y deriva la cuestión al Ministerio Público Fiscal, que consulta nuevamente a la Dirección Nacional de Migraciones; el organismo reitera su postura y, en consecuencia, se deniega otra vez la carta de ciudadanía. El interesado apela ante la Cámara Nacional Electoral, por tratarse de una cuestión con incidencia directa sobre derechos políticos.

La Cámara Electoral hace lugar al planteo y declara la nulidad del DNU 366/2025, por dos razones principales: la ausencia de necesidad y urgencia que justificara recurrir a esa vía --- pudiendo regularse la cuestión por ley ---, y la circunstancia de tratarse de una materia expresamente vedada al dictado de decretos de necesidad y urgencia.

\begin{formula}{Art. 99 inciso 3, Constitución Nacional}
``El Poder Ejecutivo no podrá en ningún caso, bajo pena de nulidad absoluta e insanable, emitir disposiciones de carácter legislativo. Solamente cuando circunstancias excepcionales hicieran imposible seguir los trámites ordinarios previstos por esta Constitución para la sanción de las leyes, y no se trate de normas que regulen materia penal, tributaria, electoral o el régimen de los partidos políticos, podrá dictar decretos por razones de necesidad y urgencia [\ldots]''

Norma central del caso: la materia electoral está expresamente vedada a los DNU.
\end{formula}

\section{La discusión sobre qué modifica exactamente el DNU}

El DNU modifica los requisitos que tenía esta persona para acceder a la ciudadanía, haciéndolos considerablemente más estrictos, requisitos que el interesado no cumplía. Esa restricción, de haberse dispuesto por ley, revestiría otra legalidad; al haberse dispuesto por DNU existe una extralimitación de competencias, en tanto el Poder Ejecutivo estaría legislando en una materia que el art.\ 99 inciso 3 no le permite.

El Poder Ejecutivo presentó, a su vez, el recurso extraordinario federal --- ya admitido y pendiente de resolución por la Corte --- argumentando que la sentencia de Cámara fue arbitraria en un doble sentido: primero, porque la nulidad debería haber tenido efecto únicamente para el caso concreto (efecto inter partes) y no derogar el decreto en su totalidad; segundo, porque el gobierno invoca razones de gestión y de seguridad para mantener la competencia migratoria en manos del Poder Ejecutivo.

\begin{note}
\textbf{Analogía señalada en clase: el decreto sobre discursos de odio.} El caso se vincula con un decreto sancionado poco tiempo antes, que habilita modificar la situación de ciudadanía o impedir el ingreso al país por incurrir en manifestaciones de discurso de odio fundadas en la nacionalidad, dictado tras un debate público posterior a un mundial de fútbol. Se señalan tres puntos de similitud: en ambos casos se recurre a un DNU; en ambos se modifica la Ley de Migraciones respecto de la posibilidad de denegar la ciudadanía o vedar el ingreso; y en ambos existe una incidencia indirecta --- pero real --- sobre los derechos políticos de las personas afectadas, lo que podría colisionar con la prohibición de dictar normas de carácter electoral por esa vía.
\end{note}

\begin{note}
\textbf{Valoración personal del docente (opinión, no dato objetivo del derecho vigente).} Se planteó que existen cuestionamientos serios respecto de este DNU: no habría una acreditación de la necesidad imperiosa y de la urgencia exigidas para recurrir a esa vía, y se advirtió que podría haber más oportunismo político --- en función del debate público existente --- que una verdadera situación excepcional, máxime cuando el Congreso se encontraba en condiciones de sesionar y el Poder Ejecutivo podría haber enviado un proyecto de ley.
\end{note}

\section{Otros aspectos discutidos: \textit{iura novit curia} y la naturaleza penal de la sanción}

Se destacan dos observaciones adicionales sobre el fallo de Cámara. En primer lugar, el fallo de Cámara pareciera resolver la inconstitucionalidad del decreto bajo el principio de \textit{iura novit curia}, dado que esa cuestión no había sido planteada expresamente por la parte actora: el tribunal inferior simplemente había denegado la ciudadanía, y recién en Cámara se hace hincapié tajante en ese aspecto.

En segundo lugar, más allá de que formalmente se la califique como infracción administrativa, la posible expulsión inmediata del país tiene un efecto fuertemente punitivo, de naturaleza prácticamente penal. En este caso es la propia Dirección Nacional de Migraciones la que decide, por sí misma, si la infracción se configura o no, sin que medie la comisión de un delito informado por otra vía. Esto plantea una doble intromisión del DNU en materias vedadas: la electoral y la penal.

Se destaca también que la persona no tenía antecedentes penales en la Argentina ni en la República Popular China, extremo que presentó debidamente traducido, lo que refuerza la desproporción de la denegatoria basada exclusivamente en la falta de registro de ingreso.

\section{Cierre del caso}

\begin{important}
\textbf{Síntesis.} Existe una cuestión federal compleja y directa: hay una colisión entre el DNU 366/2025 y el art.\ 99 inciso 3 de la Constitución Nacional, tanto por la falta de acreditación de la necesidad y urgencia como por tratarse de una materia expresamente vedada a esa herramienta normativa. El requisito de sentencia definitiva de última instancia se satisface con el pronunciamiento de la Cámara Nacional Electoral, tribunal superior en la materia, ya admitido el recurso extraordinario federal y pendiente de resolución por la Corte Suprema.
\end{important}

% ============================================================
\chapter[Caso 3 --- Movilidad jubilatoria (Badaro)]{Caso 3 --- Movilidad jubilatoria y control de constitucionalidad: el caso ``Badaro''}

\section{Los hechos}

El tercer grupo expone el caso ``Badaro c/ ANSES'', resuelto por la Corte Suprema de Justicia de la Nación en el año 2007. El señor Badaro era un jubilado cuyo haber superaba el haber mínimo. Entre los años 2002 y 2006 se otorgaron distintos aumentos a las jubilaciones, dirigidos principalmente a quienes percibían haberes mínimos --- no era el caso de Badaro, cuyo haber era algo más alto ---. Según los datos considerados por la Corte, en ese período el nivel general de precios había aumentado un 91\% y los salarios un 88\%, mientras que el haber de Badaro solamente había aumentado un 11\% en esos cuatro años.

Por ese motivo, Badaro sostuvo que su jubilación había perdido gran parte de su poder adquisitivo y reclamó judicialmente el reajuste de su haber, centrando el planteo no solo en el monto percibido al momento de jubilarse, sino en cómo debía actualizarse ese haber a lo largo del tiempo.

\begin{formula}{Artículo 14 bis, Constitución Nacional (garantía de movilidad)}
``[\ldots] jubilaciones y pensiones móviles [\ldots]''

Cláusula que consagra expresamente la movilidad de los haberes jubilatorios como garantía constitucional, núcleo del reclamo de Badaro.
\end{formula}

\section{El recorrido procesal}

El reclamo comienza en primera instancia, donde se rechaza el haber inicial. La Cámara confirma parcialmente la decisión aplicando criterios de movilidad que no le resultan favorables. El caso llega entonces a la Corte Suprema. En 2006, la Corte reconoce el problema de la movilidad, pero --- en lugar de resolverlo directamente --- le otorga un plazo al Congreso para que legisle una solución. Como las medidas adoptadas por el Congreso no resuelven el deterioro producido entre 2002 y 2006, la Corte vuelve a intervenir en el año 2007.

\section{La cuestión federal}

El problema jurídico principal consistía en determinar si el régimen de movilidad de la Ley 24.463 era compatible con la garantía de jubilaciones móviles que establece el artículo 14 bis de la Constitución. La Corte señaló que esa ley consagraba un régimen de movilidad con un nivel de protección menor al que existía al momento de su entrada en vigencia, eliminando los ajustes basados en la comparación de índices salariales.

\begin{formula}{Artículo 7, inciso 2, Ley 24.463 (Ley de Solidaridad Previsional)}
Según lo referido en la exposición, la norma estipulaba que la movilidad de los haberes jubilatorios se determinara anualmente en función de lo que dispusiera el presupuesto, facultad que en la práctica el Congreso no ejerció de forma efectiva durante el período cuestionado. Es la norma federal cuya constitucionalidad fue cuestionada frente al art.\ 14 bis de la Constitución Nacional.
\end{formula}

Se caracteriza la cuestión federal como compleja y directa: compleja, porque requiere confrontar normas de distinta jerarquía --- una ley federal frente a la Constitución Nacional ---; directa, porque el conflicto se produce de forma inmediata entre la norma legal y el texto constitucional, sin normas intermedias.

No se trataba solamente de discutir cuánto dinero debía cobrar un jubilado, sino de determinar si el Estado había cumplido con la obligación de garantizar jubilaciones móviles establecida en el artículo 14 bis. Por ese motivo la Corte debió ejercer el control de constitucionalidad sobre la Ley de Solidaridad Previsional.

\section{La decisión de la Corte y su relevancia}

El 26 de noviembre de 2007, la Corte declara la inconstitucionalidad del artículo 7, particularmente de su inciso 2, de la Ley 24.463, en cuanto se había producido una movilidad insuficiente del haber del señor Badaro. Ordena que su prestación se reajuste entre enero de 2002 y diciembre de 2006 conforme a las variaciones del índice de salarios de nivel general elaborado por el organismo oficial de estadísticas, con el pago de las diferencias correspondientes.

La movilidad jubilatoria es un derecho amparado en el artículo 14 bis, y la idea que subyace a ese derecho es que no se genere un agravio patrimonial ni se reduzca el valor del haber jubilatorio, que es de carácter alimentario y debe reemplazar al ingreso ordinario de la etapa activa. Toda movilidad debe perseguir que no se destruya el poder adquisitivo de la clase pasiva.

Lo que la Corte estableció en este caso fue fijar un estándar que asegure un piso mínimo: disminuido el haber más allá de lo que ese piso prevé, ya se genera un agravio patrimonial para la clase pasiva. Se trata de una clase de fallos de la Corte que buscan generar un estándar para que luego la clase política adopte una determinación, pero siempre sobre bases que aseguren que no exista una violación del derecho de propiedad de los jubilados y jubiladas.

\begin{note}
\textbf{Impacto posterior del precedente.} Toda la doctrina y jurisprudencia previsional posterior se apoya en este precedente, siendo base de cientos de miles de juicios de reajuste desde la Reparación Histórica de 2016. Hubo otras cláusulas de ajuste que reemplazaron leyes previsionales y que, en gestiones posteriores --- tanto anteriores como más recientes ---, también fueron declaradas inconstitucionales por motivos similares, lo que confirma que lo constitucionalmente adecuado es contar con una ley previsible que marque un horizonte claro de actualización periódica de los haberes.
\end{note}

\section{Cierre del caso}

\begin{important}
\textbf{Síntesis.} ``Badaro'' es un caso especialmente útil para la consigna de la clase porque reúne todos los elementos: una garantía constitucional expresa --- la movilidad jubilatoria del art.\ 14 bis ---, una norma federal cuya constitucionalidad se cuestiona, una cuestión federal compleja y directa, y la intervención de la Corte Suprema por vía del recurso extraordinario federal.
\end{important}

% ============================================================
\chapter[Caso 4 --- Poder tributario municipal]{Caso 4 --- Poder tributario municipal, coparticipación y cláusula comercial}

\section{Los hechos}

El cuarto grupo expone un caso vinculado a la ``tasa vial municipal'' sancionada en 2024. La controversia surge porque numerosos municipios de distintas provincias aprobaron ordenanzas tributarias que aplicaban un recargo directo sobre los combustibles --- entre el 1,5\% y el 2\% por litro de nafta y de gasoil cargado ---. La medida generó reclamos de cámaras empresariales, de la cámara que nuclea a las estaciones de servicio y de la Secretaría de Energía de la Nación, por considerarse gravámenes ilegítimos. Se cita como ejemplo el reclamo iniciado por un ex diputado en la ciudad de Mar del Plata.

\section{El problema constitucional de fondo}

La falla principal de la norma es la distorsión de la figura de la tasa, que exige una contraprestación individualizada y directa de un servicio hacia el contribuyente. Si una persona simplemente está de paso y carga combustible en una estación de servicio, no existe una contraprestación real hacia ese usuario; según el criterio expuesto en clase, la norma municipal está en realidad creando un impuesto encubierto.

\begin{formula}{Artículo 9, Ley 23.548 --- Ley de Coparticipación Federal de Impuestos}
Según lo referido en la exposición, la norma prohíbe a las provincias y a sus municipios aplicar gravámenes locales análogos a los tributos nacionales coparticipados --- el llamado ``principio de no analogía tributaria'' ---, y el impuesto sobre los combustibles ya se encuentra gravado a nivel nacional. Es el fundamento central para sostener que la tasa vial configura una doble imposición vedada por la ley-convenio de coparticipación.
\end{formula}

\begin{formula}{Artículo 75 inciso 13, Constitución Nacional --- Cláusula comercial}
``Corresponde al Congreso [\ldots] reglar el comercio con las naciones extranjeras, y de las provincias entre sí.''

La regulación del comercio interprovincial es atribución exclusiva y excluyente del Congreso de la Nación; un recargo local sobre el combustible encarece de forma directa el transporte de cargas interjurisdiccional.
\end{formula}

\begin{formula}{Artículos 9, 10 y 11, Constitución Nacional --- Prohibición de aduanas interiores}
Según el núcleo de estas cláusulas, en todo el territorio nacional no puede haber más aduanas que las nacionales, y la circulación de bienes de producción o fabricación nacional entre provincias debe ser libre de derechos de tránsito. El grupo sostiene que tasas de este tipo crean ``fronteras fiscales invisibles'' entre provincias, violando este principio.
\end{formula}

También está en juego el artículo 17, el derecho de propiedad, por la confiscatoriedad que puede llegar a generar el tributo, cuya alícuota se ubicaba entre el 1,5\% y el 2\% por litro.

\section{El precedente ``Municipalidad de Quilmes'' y la discusión sobre tasas e impuestos}

El precedente jurisprudencial más importante en esta materia es ``Municipalidad de Quilmes'', sobre el que existe una larga discusión doctrinaria. Los impuestos son exacciones patrimoniales que el Estado percibe sin una contraprestación directa, mientras que la tasa, a diferencia del impuesto, retribuye una contraprestación concreta y por lo tanto debe guardar proporcionalidad con el servicio efectivamente prestado. Que la tasa se exprese, en general, como un porcentaje de una operación está muy discutido, porque ello la emparenta con la figura del impuesto. La Corte, en ese precedente, aceptó que hubiera tasas percibidas como un porcentaje de un acto o de una transacción, y sostuvo que esa expresión porcentual, en sí misma, no imposibilita la legalidad de la tasa --- aunque el punto sigue siendo muy cuestionado.

\begin{note}
\textbf{El caso ``Rafo'' (provincia de Córdoba).} Como ejemplo paradigmático de litigio sobre tasas municipales, se menciona el caso ``Rafo'', que recorrió todas las instancias en la provincia de Córdoba: comenzó en el fuero contencioso administrativo, pasó por la Cámara Contenciosa y llegó al Tribunal Superior de Justicia de Córdoba en apelación y casación, instancia en la que se rechazó el recurso extraordinario federal. Se trata de un caso distinto del de Mar del Plata expuesto por el grupo, aunque comparable, ya que fueron muchos los municipios de distintas provincias que sostuvieron tasas de este tipo.
\end{note}

\section{Caracterización de la cuestión federal}

La cuestión constitucional se caracteriza, principalmente, porque va en contra de la ley de coparticipación, por el principio de no analogía: un impuesto no puede ser igual a otro, y en este caso, con el tema de los combustibles, el impuesto ya lo percibe la Nación. Eso configuraría una doble imposición. Se trata, entonces, de una cuestión federal compleja indirecta, porque la norma inferior --- la ordenanza municipal --- resulta contraria a otra norma también infraconstitucional, la ley de coparticipación.

Se agrega que, en paralelo, existe también una posible interferencia directa con la cláusula de comercio interjurisdiccional del artículo 75 inciso 13 de la Constitución, lo que habilitaría a plantear, además, una cuestión federal compleja directa por colisión inmediata con el propio texto constitucional.

\section{Vías procesales disponibles}

Las vías procesales para el tema de las tasas pueden ser la vía ordinaria, con todas las instancias, hasta llegar al recurso extraordinario federal, o bien una acción declarativa de certeza, o una medida de no innovar; en general, el amparo resulta menos idóneo que la acción declarativa de certeza para este tipo de tributos, porque esta última se promueve antes de que se produzca la imposición efectiva.

La acción declarativa de certeza se utiliza mucho en materia tributaria porque, antes de que empiece a percibirse el tributo, permite plantear al juez que existe una incertidumbre sobre si corresponde pagarlo o no, por dos motivos posibles: entender si la norma es aplicable --- es decir, si se es sujeto imponible --- o entender si la norma es legítima y, por lo tanto, obligatoria. Por cualquiera de esas dos vías, la acción declarativa de certeza es una herramienta idónea para impugnar preventivamente ese tipo de tasa.

\section{Cierre del caso}

\begin{important}
\textbf{Síntesis.} La tasa vial municipal admite dos vías de cuestión federal: una compleja indirecta, por colisión con la ley de coparticipación --- norma infraconstitucional ---, y otra compleja directa, por posible violación de la cláusula comercial y de la prohibición de aduanas interiores previstas en el propio texto constitucional. Ambas vías son perfectamente atendibles como argumento constitucional y muestran cómo un mismo hecho puede dar lugar a más de un tipo de cuestión federal.
\end{important}

% ============================================================
\chapter[Caso 5 --- Autonomía universitaria]{Caso 5 --- Autonomía universitaria, insistencia legislativa y medida cautelar}

El quinto grupo aclara, antes de comenzar, que el caso elegido aún no cuenta con una sentencia definitiva, pero lo considera valioso por involucrar un recurso extraordinario federal ya utilizado por el gobierno nacional durante el trámite de una medida cautelar.

\section{Los hechos y el contexto normativo}

En 2024 el Congreso sanciona la Ley 27.757, de financiamiento universitario, que el Poder Ejecutivo veta en su totalidad. El Congreso intenta insistir en la sanción, pero no logra reunir la mayoría necesaria. En 2025 se sanciona una nueva ley, la Ley 27.795 --- en octubre de ese año ---, que establece la actualización del presupuesto universitario, de los salarios docentes y no docentes y de las becas estudiantiles. El Poder Ejecutivo vuelve a vetarla, pero en esta oportunidad el Congreso logra la insistencia legislativa y la ley queda sancionada.

\begin{formula}{Artículo 83, Constitución Nacional --- Insistencia legislativa}
``Desechado en el todo o en parte un proyecto por el Poder Ejecutivo, vuelve con sus objeciones a la Cámara de su origen; ésta lo discute de nuevo, y si lo confirma por mayoría de dos tercios de votos, pasa otra vez a la Cámara de revisión. Si ambas Cámaras lo sancionan por igual mayoría, el proyecto es ley y pasa al Poder Ejecutivo para su promulgación.''

Mecanismo que le permitió al Congreso sancionar definitivamente la Ley 27.795 pese al veto presidencial.
\end{formula}

El Poder Ejecutivo dicta entonces el Decreto 759/2025, mediante el cual promulga la ley, pero condiciona la ejecución de sus artículos 5 y 6 --- referidos a la recomposición y actualización de los salarios docentes y no docentes --- a que se determine previamente la fuente de financiamiento correspondiente.

\begin{formula}{Artículos 5 y 6, Ley 27.795}
Según lo referido en la exposición, estos artículos regulan la recomposición y actualización de los salarios docentes y no docentes universitarios y de las becas estudiantiles --- entre ellas, las Becas Manuel Belgrano ---, cuya ejecución quedó condicionada por el Decreto 759/2025. Constituyen el núcleo normativo del conflicto entre el Poder Ejecutivo y el Consejo Interuniversitario Nacional.
\end{formula}

\section{El amparo colectivo y la medida cautelar}

Ante la suspensión de hecho de esos artículos, el Consejo Interuniversitario Nacional inicia una acción de amparo colectivo, solicitando la declaración de inconstitucionalidad del decreto condicionante, la aplicación directa de la ley y, como medida cautelar, que se exija al Poder Ejecutivo el cumplimiento inmediato de los artículos 5 y 6.

En diciembre de 2025, la primera instancia del fuero contencioso administrativo hace lugar a la cautelar, entendiendo que la insistencia legislativa del Congreso --- en los términos del artículo 83 de la Constitución --- tiene prevalencia sobre la voluntad posterior del Ejecutivo de condicionar la norma. El Poder Ejecutivo apela esa decisión.

En marzo de 2026, la Cámara confirma el fallo de primera instancia, sosteniendo que el Poder Ejecutivo no logró desvirtuar la verosimilitud del derecho invocado por el Consejo Interuniversitario Nacional, y que no puede ampararse en un decreto anterior cuando el Congreso sanciona posteriormente una ley que dispone precisamente lo contrario. La Cámara entiende, además, que el incumplimiento de la ley genera un perjuicio concreto sobre los salarios docentes y no docentes, con repercusión directa sobre el derecho de enseñar y aprender.

\section{El intento de levantamiento de la cautelar}

Meses después, el Poder Ejecutivo intenta levantar la medida cautelar argumentando que las circunstancias habían cambiado. Para sostener ese argumento invoca un acta firmada el 10 de junio de 2026 entre la Subsecretaría de Políticas Universitarias y distintos sectores gremiales, en la que se acordaron, entre otras medidas, un aumento salarial, una actualización de los gastos de funcionamiento de las universidades, un incremento de las Becas Manuel Belgrano y un refuerzo presupuestario para los hospitales universitarios.

\begin{note}
\textbf{La postura del Estado frente a la cautelar.} El Estado sostuvo que, habiendo adoptado ya estas medidas, la situación que había dado origen a la cautelar --- el incumplimiento de los arts.\ 5 y 6 de la Ley 27.795 --- se encontraba momentáneamente resuelta, por lo que ya no existía necesidad de mantenerla vigente.
\end{note}

Sin embargo, el 1.\textsuperscript{o} de septiembre de 2026 el juez de la causa rechaza el pedido de levantamiento, por considerar que el Estado no había demostrado de manera suficiente que esas medidas hubieran solucionado efectivamente el problema; entendió que persistía el perjuicio y que seguía vigente el peligro en la demora que había justificado originalmente la cautelar. El juez adoptó, además, una postura de seguimiento del cumplimiento de la medida: dictó una resolución ``para mejor proveer'' y otorgó al Estado un plazo de tres días para presentar documentación e información concreta sobre el estado de cumplimiento de los compromisos asumidos y de la propia cautelar.

\begin{note}
\textbf{El precedente ``Mendoza'' como antecedente de seguimiento judicial.} Se menciona como precedente el caso ``Mendoza'', en el que la Corte Suprema también estableció un sistema de seguimiento y control judicial para garantizar el cumplimiento efectivo de sus decisiones, más allá de la sentencia formal.
\end{note}

\section{El recurso extraordinario federal y su admisibilidad: el punto crítico}

La valoración crítica del grupo sobre la procedencia del recurso utilizado por el Poder Ejecutivo sostiene que, si bien existe un caso de cuestión federal y un recurso extraordinario federal efectivamente utilizado por el Poder Ejecutivo, dicho recurso estaría mal utilizado, porque falta el requisito fundamental de la sentencia definitiva: se trata de un recurso aplicado sobre una medida cautelar dentro de un amparo colectivo aún en trámite. La doctrina admite casos de excepción por ``gravedad institucional'', que es lo que invoca el gobierno nacional, pero la Corte, en definitiva, rechaza ese planteo.

La cuestión federal remite, en este caso, a la autonomía de las universidades nacionales, consagrada en el artículo 75 inciso 19 de la Constitución Nacional, y en definitiva al cumplimiento de la propia ley.

\begin{formula}{Artículo 75 inciso 19, Constitución Nacional --- Autonomía universitaria}
``Corresponde al Congreso: [\ldots] Sancionar leyes de organización y de base de la educación que consoliden la unidad nacional [\ldots] y que garanticen los principios de gratuidad y equidad de la educación pública estatal y la autonomía y autarquía de las universidades nacionales.''

Cláusula del progreso. Fundamento constitucional de la autonomía y autarquía universitaria invocado por el Consejo Interuniversitario Nacional.
\end{formula}

Se agrega también que la Cámara hizo hincapié en el artículo 83 de la Constitución en cuanto regula la insistencia del Congreso: si el Congreso insiste tanto en una ley, esa ley debe tener un cumplimiento efectivo por parte del Poder Ejecutivo.

\section{Cierre del caso}

\begin{important}
\textbf{Síntesis.} El caso presenta una cuestión federal compleja y directa --- colisión entre un decreto del Poder Ejecutivo y la Constitución Nacional en materia de autonomía universitaria e insistencia legislativa ---, pero con un problema de admisibilidad procesal: el recurso extraordinario federal fue interpuesto sobre una medida cautelar, sin sentencia definitiva, lo que en principio impide su procedencia salvo que se acredite una verdadera gravedad institucional, extremo que la Corte no reconoció en este caso.
\end{important}

% ============================================================
\chapter[Unidad 6 --- Cuadro comparativo]{Unidad 6 --- Cuadro comparativo de los cinco casos}

El siguiente cuadro sintetiza, para cada caso expuesto, la norma inferior cuestionada, la norma constitucional o de mayor jerarquía en tensión, el tipo de cuestión federal identificada y el estado procesal del recurso extraordinario federal al momento de la clase.

\begin{table}[H]
\centering
\footnotesize
\begin{tabularx}{\textwidth}{
    >{\raggedright\arraybackslash}p{0.14\textwidth}
    >{\raggedright\arraybackslash}p{0.19\textwidth}
    >{\raggedright\arraybackslash}p{0.24\textwidth}
    >{\raggedright\arraybackslash}p{0.14\textwidth}
    >{\raggedright\arraybackslash}X
}
\toprule
\textbf{Caso} & \textbf{Norma inferior cuestionada} & \textbf{Norma constitucional en juego} & \textbf{Tipo de cuestión federal} & \textbf{Estado del REF} \\
\midrule
1. Vivienda digna (CABA) & Interpretación restrictiva de la Ley 4036 (CABA) & Arts.\ 14 bis y 16 CN; arts.\ 17 y 18 CCABA; CEDAW; PIDESC & Compleja, directa & Presentado (ago.\ 2026), pendiente ante el TSJ \\
2. Ciudadanía y migraciones & DNU 366/2025 & Art.\ 99 inc.\ 3 CN (materia electoral vedada) & Compleja, directa & Admitido, pendiente ante la Corte \\
3. Movilidad jubilatoria (Badaro) & Art.\ 7 inc.\ 2, Ley 24.463 & Art.\ 14 bis CN (movilidad) & Compleja, directa & Resuelto por la Corte (2007) \\
4. Tasa vial municipal & Ordenanzas municipales tributarias & Ley 23.548 (coparticipación) y arts.\ 9, 10, 11, 17 y 75 inc.\ 13 CN & Compleja, indirecta y directa & Múltiples reclamos en trámite en distintas provincias \\
5. Financiamiento universitario & Decreto 759/2025 & Art.\ 83 CN (insistencia) y art.\ 75 inc.\ 19 CN (autonomía) & Compleja, directa & Interpuesto sobre cautelar; problema de admisibilidad \\
\bottomrule
\end{tabularx}
\caption{Cuadro comparativo de los cinco casos expuestos en clase.}
\end{table}

% ============================================================
\chapter[Unidad 7 --- Resumen Cornell]{Unidad 7 --- Resumen Cornell y síntesis final de la clase}

Herramienta de repaso activo: la columna izquierda contiene preguntas disparadoras y conceptos clave; la columna derecha, la respuesta o el desarrollo correspondiente.

\begin{longtable}{
    >{\raggedright\arraybackslash}p{0.30\textwidth}
    >{\raggedright\arraybackslash}p{0.60\textwidth}
}
\caption{Cuadro Cornell de repaso de la clase.} \\
\toprule
\textbf{Preguntas / palabras clave} & \textbf{Notas / respuestas} \\
\midrule
\endfirsthead

\toprule
\textbf{Preguntas / palabras clave} & \textbf{Notas / respuestas} \\
\midrule
\endhead

\midrule
\multicolumn{2}{r}{\textit{Continúa en la página siguiente}} \\
\endfoot

\bottomrule
\endlastfoot


\textbf{¿Qué es una ``cuestión federal''?} &
Es el conflicto entre una norma de rango inferior (ordenanza, decreto, ley local o federal) y la Constitución Nacional, un tratado o una ley federal, cuya resolución habilita, eventualmente, el recurso extraordinario federal ante la Corte Suprema. \\

\textbf{¿Cuáles son los tipos de cuestión federal vistos en clase?} &
Compleja directa (colisión inmediata entre una norma y la Constitución, como en los casos de vivienda, ciudadanía, jubilaciones y universidades) y compleja indirecta (colisión entre dos normas infraconstitucionales de distinta jerarquía, como la ordenanza municipal frente a la ley de coparticipación en el caso de la tasa vial). \\

\textbf{¿Qué requisitos exige el recurso extraordinario federal?} &
Existencia de un caso judicial, sentencia definitiva dictada por el superior tribunal de la causa, agotamiento de las instancias ordinarias previas y relación directa entre la cuestión federal y la resolución del caso (art.\ 14, Ley 48). \\

\textbf{Caso 1 --- ¿Por qué no existe un orden de prelación para acceder a la vivienda?} &
Porque ni los arts.\ 14 bis y 16 CN, ni la Ley 4036, ni los tratados internacionales (CEDAW, PIDESC) establecen un ranking de beneficiarios; por criterio hermenéutico, donde la norma no distingue, no cabe distinguir para denegar un derecho. \\

\textbf{Caso 2 --- ¿Por qué se declaró nulo el DNU 366/2025?} &
Porque no se acreditó necesidad y urgencia y porque transfería a un órgano administrativo una materia con incidencia electoral, expresamente vedada al Poder Ejecutivo por el art.\ 99 inciso 3 de la Constitución. \\

\textbf{Caso 3 --- ¿Qué estableció la Corte en ``Badaro''?} &
Que la movilidad jubilatoria del art.\ 14 bis no puede quedar librada a la discrecionalidad estatal; fijó un piso mínimo de protección y ordenó reajustar el haber conforme a índices salariales objetivos, sentando un precedente central en materia previsional. \\

\textbf{Caso 4 --- ¿Por qué la tasa vial es cuestionable?} &
Porque distorsiona la figura de la tasa --- que exige contraprestación directa --- convirtiéndola en un impuesto encubierto; viola el principio de no analogía tributaria (Ley de Coparticipación) y afecta la cláusula comercial y la prohibición de aduanas interiores (arts.\ 9, 10, 11 y 75 inc.\ 13 CN). \\

\textbf{Caso 5 --- ¿Por qué el recurso extraordinario federal estaba ``mal utilizado''?} &
Porque se interpuso contra una medida cautelar dentro de un amparo colectivo aún en trámite, sin sentencia definitiva de última instancia --- requisito esencial del recurso, salvo excepcional gravedad institucional, no reconocida por la Corte en este caso. \\

\textbf{¿Qué diferencia una tasa de un impuesto?} &
El impuesto es una exacción patrimonial sin contraprestación directa; la tasa, en cambio, retribuye una contraprestación estatal concreta y debe guardar proporcionalidad con ese servicio (precedente ``Municipalidad de Quilmes''). \\

\textbf{¿Qué es la ``insistencia legislativa'' del art.\ 83 CN?} &
El mecanismo por el cual el Congreso puede sancionar definitivamente una ley pese al veto presidencial, si ambas Cámaras la confirman por mayoría de dos tercios de votos; una vez lograda, el Ejecutivo no puede condicionar su ejecución por decreto. \\

\midrule

\multicolumn{2}{p{0.90\textwidth}}{
\textbf{Resumen:} A lo largo de la clase, cinco casos muy distintos entre sí --- una mujer en situación de vulnerabilidad que reclama vivienda, un comerciante extranjero al que le niegan la ciudadanía, un jubilado cuyo haber pierde poder adquisitivo, automovilistas que pagan una tasa municipal sobre el combustible y un sistema universitario que reclama el cumplimiento de su presupuesto --- permitieron trabajar sobre una misma herramienta conceptual: la cuestión federal y el recurso extraordinario federal como vía de acceso a la Corte Suprema de Justicia de la Nación. En todos los casos aparece el mismo patrón: una norma de rango inferior entra en tensión con una norma de rango superior, y esa tensión solo puede resolverse en última instancia por la Corte Suprema, siempre que se cumplan los requisitos formales del recurso: sentencia definitiva del superior tribunal de la causa, agotamiento de las vías ordinarias y relación directa entre la cuestión federal planteada y la solución del pleito --- requisito que, como se vio en el caso de las universidades, no siempre se encuentra satisfecho ---. El control de constitucionalidad no es, en definitiva, un ejercicio abstracto: es la herramienta con la que el sistema jurídico argentino corrige, caso por caso, los desbordes de competencia de los poderes públicos frente a los derechos y garantías que la Constitución promete a cada persona.
} \\

\end{longtable}

\end{document}
